\section{Limitations}
\label{sec:limitations}

\begin{enumerate}[leftmargin=1.4em,itemsep=2pt]
\item \textbf{Compute.} All results come from one $8$\,GB laptop GPU and roughly
$30$ GPU-hours. This single constraint explains most of the gaps below. We state
it once rather than repeating it as a caveat on every claim.
\item \textbf{One detector family, two depths on the clean split.} YOLOv8s
(five seeds) and YOLOv8n (three seeds). YOLOv8m was trained only on the old leaky
split and is therefore not comparable; we exclude it from the size sweep rather
than presenting a contaminated point.
\item \textbf{Small-model seed count, and shared components.} The \texttt{yolov8n}
row now carries an interval over three seeds, but its test mAP sits inside the
\texttt{yolov8s} interval, so the two are not distinguishable on that metric. The
five two-stage seeds are also not five independent systems. All five are
detector-aligned at evaluation time, each scoring its own stage-1 detections, but
seeds 0--2 train their classifier on a single crop set generated once from the
seed-0 detector, while seeds 3--4 generate their own crops. The stage-2 recipe is
held fixed throughout, and the stage-2 classifiers themselves are separately
trained (distinct checkpoints per seed). The effective replication is closer to
four.
\item \textbf{Unresolved checks.} We did not sweep recall targets and FP budgets
exhaustively, did not threshold-match the leaky and clean models, and did not
retrain after the label correction. The leak count rests on the operational
filter plus an embedding-neighbour check that agreed on all $99$ survivors but
cannot exclude a single missed near-duplicate; one candidate pair
(\texttt{ds1\_0054}/\texttt{ds1\_0111}) is reported in the run log as unconfirmed.
\item \textbf{Five negative test images.} Precision and FP/image are
high-variance; the source-cluster intervals span an order of magnitude. Point
estimates in tables are accompanied by intervals wherever the interval is
meaningful.
\item \textbf{Label correction without retraining.} Two test boxes were removed
from ground truth. We report both versions, but the model was not retrained on
corrected labels because no training-label errors were found.
\item \textbf{No grouped cross-validation.} A single frozen split was used, as
required by the protocol. Group-aware $k$-fold would give a better estimate of
split variance and was not affordable.
\item \textbf{No hard-negative re-mining on unseen data.} The stage-2 negatives
come from detector false positives on the training split, which is the correct
procedure but limits how novel the negatives are.
\item \textbf{Synthetic composites.} The paste-and-remove experiment uses
feathered, scale-mismatched composites. It is a probe, not a counterfactual.
\item \textbf{No on-device characterisation.} All latency numbers in
Table~\ref{tab:size} are PC GPU FP32 at $640\times640$. No edge measurement is
reported here, so we make no claim about Jetson throughput or energy. The
accompanying AeroShield deployment project does serve this same single-class
model on a Jetson Nano via TensorRT and exposes a per-run timing summary, so the
gap is an instrumentation task rather than a change of model.
\item \textbf{Source confound.} \texttt{Jan} and \texttt{May} are colour
photographs with marker pins, despite \texttt{lwir} in the filename. Marker
shortcuts cannot be fully excluded, and the source dominates aggregate false
alarms.
\item \textbf{No field validation and no buried-object claim.} This is a
benchmark on visible surface indicators. It says nothing about burial depth,
soil type, or detection probability in real clearance conditions.
\end{enumerate}
